\documentclass[12pt]{article}
\usepackage{report}

\begin{document}

\reporttitle{CS340400 Compiler Design - Homework 1 Implementation Report}

\section{Overview}

This assignment implements a scanner for a subset of C with flex. The scanner reads source code,
prints every token with its category, and after each line prints the line number together with
the original text of that line. Two directives, \code{#pragma source on|off} and
\code{#pragma token on|off}, turn each kind of output on or off.

\begin{center}
\begin{tabularx}{0.92\linewidth}{@{}lL@{}}
\toprule
Files & \file{scanner.l}, \file{makefile} \\
Tests & 8 cases in \file{tests/}: one per grading item plus a comprehensive case --- \textbf{8/8 pass} \\
Build & GNU gcc 13.3, \textbf{0 warnings} (also with \code{-Wall}) \\
Usage & \code{make scanner && ./scanner < input.c}; \code{make test} runs all cases \\
\bottomrule
\end{tabularx}
\end{center}

Pragmas, escape sequences, patterns inside comments and a last line without a newline are the places
where a straightforward scanner goes wrong; Section~\ref{sec:diff} summarises how each one is handled.

\section{Output Format}

The specification asks for two kinds of output, interleaved in input order:

\begin{itemize}
  \item every token as \code{#<category>:<text>}, e.g.\ \code{#key:int}, \code{#id:a}, \code{#op:=};
  \item at the end of every line, \code{<line number>:<line text>}, e.g.\ \code{2:char a = 'i';}.
\end{itemize}

Comments and pragmas produce no tokens, but the lines that contain them are still printed.
\code{#pragma token off} silences only the tokens, and \code{#pragma source off} silences only the
lines.

\begin{lstlisting}[language={},numbers=none,title={\codetitle{input: int a = 1; // x}}]
#key:int
#id:a
#op:=
#integer:1
#punc:;
1:int a = 1; // x
\end{lstlisting}

\section{Structure}

The scanner is a single flex file with the usual three parts:

\begin{center}
\begin{tabularx}{0.92\linewidth}{@{}lL@{}}
\toprule
Part & Contents \\
\midrule
Definitions \code{\%{ ... \%}} & line buffer, pragma switches, output helpers, \code{YY_USER_ACTION} \\
Definitions (patterns) & \code{\%option}, \code{\%x COMMENT}, named patterns (\code{ESCAPE}, \code{STRING}, \code{INTEGER}, \ldots) \\
Rules & comments, pragmas, keywords, macros, identifiers, literals, operators, punctuation, whitespace, EOF, error \\
User code & \code{main()} calls \code{yylex()} once \\
\bottomrule
\end{tabularx}
\end{center}

The \emph{order} of the rules matters. When several rules match, flex takes the longest match, and
breaks ties in favour of the rule written first. Consequently:

\begin{itemize}
  \item \code{for} matches both the keyword rule and \code{{IDENTIFIER}} with the same length, so the
        keyword rule, which comes first, wins;
  \item \code{format} matches 6 characters as an identifier but only 3 as a keyword, so it is an identifier;
  \item macros such as \code{NULL} and \code{INT_MAX} are placed before \code{{IDENTIFIER}} for the same reason,
        while \code{NULLX} is longer and therefore an identifier.
\end{itemize}

\subsection{Output helpers}

\snippet[language=flex]{../scanner.l}{static void append_line(}{3@end}

\begin{itemize}
  \item \code{append_line} appends the matched text to \code{line_buf}; anything beyond 4096 bytes is
        cut off instead of overflowing the buffer.
  \item \code{end_line} is called by the newline rule. It prints the line if \code{source_on} is set,
        then clears the buffer and increments the line number.
  \item \code{emit} is called by every token rule and prints the token if \code{token_on} is set.
\end{itemize}

A common alternative is to give each token an enum value and decide its category in \code{main} with
range checks such as \code{token >= 30 && token <= 50}. Calling \code{emit("key")} directly in the
rule keeps the category next to the pattern, needs no token table, and cannot break when an enum is
reordered.

\section{Implementation Details}

\subsection{Collecting the line text with \texttt{YY\_USER\_ACTION}}

\snippet[language=flex]{../scanner.l}{define YY_USER_ACTION}{+1}

flex expands \code{YY_USER_ACTION} before the action of \emph{every} rule. Because of this, whatever a
rule matches --- a keyword, whitespace, a comment or a pragma --- is appended to \code{line_buf}
automatically, and no rule has to do it by hand. The newline is the only exception: it marks the
end of the line and is handled by the \code{\n} rule, which calls \code{end_line()}.

This guarantees that the printed line is always complete, including the parts that produce no
tokens, such as comments, pragmas and whitespace.

\subsection{Pragmas}

\snippet[language=flex]{../scanner.l}{"#pragma"}{+2}
\snippet[language=flex]{../scanner.l}{{PRAGMA}source}{+4}

Two details make the pragma rules work:

\begin{enumerate}
  \item The pattern is anchored with \code{^}, so a pragma must start a line (indentation is allowed).
  \item The pattern does \emph{not} include the newline. If the pragma rule consumed \code{\n}, \code{end_line()}
        would not run for that line: the pragma line would merge with the next one and every following
        line number would be off by one. Leaving the newline to the \code{\n} rule keeps both the
        numbering and the line text correct.
\end{enumerate}

The switches also take effect at the right moment. \code{#pragma source off} clears
\code{source_on} \emph{before} its own newline is scanned, so the directive line itself is not printed,
while \code{#pragma source on} prints itself. Test \file{07_pragma} checks exactly this:

\begin{lstlisting}[language={},numbers=none,title={\codetitle{tests/07\_pragma: input (left) and expected output, abridged (right)}}]
1  #pragma source on        1:#pragma source on
2  char a = 'i';            #key:char ... #punc:;   2:char a = 'i';
3  #pragma<TAB>source  off  (not printed)
4  int b;                   #key:int #id:b #punc:;  (line text not printed)
5  #pragma source on        5:#pragma source on
6  #pragma token off        6:#pragma token off
7  int c;                   7:int c;                (no tokens)
8  #pragma token on         8:#pragma token on
9  int d;                   #key:int #id:d #punc:;  9:int d;
\end{lstlisting}

\subsection{Comments: an exclusive start condition}

\snippet[language=flex]{../scanner.l}{"//"}{+6}

Block comments use the start condition \code{COMMENT}, declared with \code{\%x} as \emph{exclusive}: while
inside it, only rules prefixed with \code{<COMMENT>} are active. With \code{\%s} (inclusive), the other rules
would stay active inside comments. For \code{/*"abc*/}, the rule for an unterminated string would
match a longer prefix than the comment rule, and the scanner would report an error.

The body of a comment is split into three rules, none of which crosses a newline:

\begin{center}
\begin{tabularx}{0.92\linewidth}{@{}lL@{}}
\toprule
Rule & Matches \\
\midrule
\code{[^*\n]+} & ordinary text without \code{*} \\
\code{"*"+[^*/\n]*} & one or more \code{*} not followed by \code{/}, e.g.\ \code{**x} \\
\code{"*"+"/"} & the end of the comment, including \code{**/} \\
\code{\n} & a line break inside the comment: call \code{end_line()} as usual \\
\bottomrule
\end{tabularx}
\end{center}

A single pattern for the whole comment would swallow the newlines, and the lines of a multi-line
comment could no longer be printed one by one. With the split rules, \code{YY_USER_ACTION} still
collects the comment text and each newline ends its line normally.

\subsection{Strings and character constants}

\snippet[language=flex]{../scanner.l}{ESCAPE}{+3}

A string body is any number of three alternatives: a valid escape sequence, a backslash followed by
any character (\code{\\.}, for escapes the specification does not list), or any character other than
\code{"}, \texttt{\textbackslash} and newline.

Because \code{\"} is consumed as one escape, \code{"a\"b"} is a single string and does not end at the
inner quote. A simpler pattern such as \code{\042[^\042\n]*\042} would stop there.
\code{"a\nb"} is likewise one token, \code{#string:"a\nb"}, quotes included, not three tokens split at
the escape. A character constant uses the same alternatives but
allows exactly \emph{one} of them between the quotes.

\subsection{Numbers}

\snippet[language=flex]{../scanner.l}{INTEGER}{+2}

The specification treats a sign as part of the number: \code{-0} is one integer, while \code{- 0} is an
operator followed by an integer. Floating-point numbers accept the forms \code{1.}, \code{.5} and
\code{1.5}. Because flex prefers the longest match, \code{a-1} is scanned as \code{#id:a} followed by
\code{#integer:-1}; Section~\ref{sec:assume} discusses this.

\subsection{End of file}

\snippet[language=flex]{../scanner.l}{<<EOF>>}{+5}

A line is printed when its newline is scanned, so without extra handling a last line that has no
trailing newline would be lost. The \code{<<EOF>>} rule is active in both start conditions and prints the
buffer if it still holds text, which also covers a comment that runs to the end of the file.

\subsection{Errors and compiler warnings}

\snippet[language=flex]{../scanner.l}{lex_error(); }{+1}

\begin{itemize}
  \item The catch-all rule \code{.} comes last. Any character that no other rule accepts prints
        \code{Error at line N: <char>} and exits with status 1.
  \item \code{\%option noyywrap} removes the need to link with \code{-lfl}; \code{nounput} and
        \code{noinput} stop flex from generating the unused functions \code{yyunput} and \code{input},
        which would trigger unused-function warnings under \code{-Wall}.
  \item The file includes \code{<stdlib.h>} for \code{exit()}; without it GNU gcc reports an implicit
        declaration, and the grading policy deducts 20 points for any warning.
\end{itemize}

\section{Testing}

\code{make test} runs \file{./scanner} on every \file{tests/*.c} and compares the output with the
matching \file{.expected} file; a failing case writes its diff to \file{tests/*.diff}.

\begin{center}
\begin{tabularx}{\linewidth}{@{}lP{4.4cm}L@{}}
\toprule
Test & Grading item (points) & Focus \\
\midrule
\file{01_keyword_id} & Keywords \& Identifiers (10) & every keyword; \code{Int}, \code{_abc}, \code{x_1} are identifiers \\
\file{02_macro} & Macros (10) & all 9 macros; \code{NULLX} is longer than \code{NULL}, so it is an identifier \\
\file{03_op_punc} & Operators \& Punctuations (15) & two-character operators win over one-character ones \\
\file{04_number_char} & Integers, Floats, Characters (15) & signs (\code{-0} vs.\ \code{- 0}), \code{5.}, \code{.5}, escapes \\
\file{05_string} & Strings (15) & \code{\"} and \code{\n} inside one token; \code{//} and \code{/*} inside strings \\
\file{06_comment} & Comments (15) & multi-line comments, \code{**}, \code{/**/}, a quote inside a comment \\
\file{07_pragma} & Pragma Directives (10) & all four pragmas, tabs and extra spaces, when the switch applies \\
\file{08_comprehensive} & Comprehensive (10) & everything together; the last line has no newline \\
\bottomrule
\end{tabularx}
\end{center}

All 8 cases pass.

\section{Pitfalls and How They Are Handled}
\label{sec:diff}

\begin{center}
\begin{tabularx}{\linewidth}{@{}rLL@{}}
\toprule
\# & Pitfall & Handling \\
\midrule
1 & A pragma pattern that consumes the newline shifts every later line number & No \code{\n} in the pattern; anchor with \code{^} \\
2 & A naive string pattern ends early at \code{\"} & \code{STRING} repeats ``escape or ordinary character'' \\
3 & Escape sequences can split one string into several tokens & The whole string is one token \\
4 & With an inclusive start condition, a quote inside a comment causes an error & \code{\%x COMMENT} \\
5 & A last line without a newline is never printed & \code{<<EOF>>} rule prints the remaining buffer \\
6 & A missing \code{<stdlib.h>} or \code{while (token = yylex())} warns under \code{-Wall} & Include it; \code{main} calls \code{yylex()} once \\
7 & Deciding token categories by enum ranges is fragile & Each rule calls \code{emit("<category>")} \\
\bottomrule
\end{tabularx}
\end{center}

\section{Assumptions}
\label{sec:assume}

The course's reference implementation, \code{golden_scanner}, is no longer available, so the following
behaviour is inferred from the specification and could not be checked against it:

\begin{itemize}
  \item A string is printed as a single token (Section 4.4).
  \item \code{a-1} is scanned as \code{#id:a}, \code{#integer:-1}. The specification only contrasts
        \code{- 0} with \code{-0}; flex's longest-match rule produces this result.
  \item A last line without a trailing newline is still printed with its line number.
\end{itemize}

\end{document}
